\section{Terminal trace estimate and quantitative closure}
\label{sec:closure}

We now insert the trace estimate from \cref{sec:terminal} into the last sum
of the descent.  This section also retains both branches of $\Delta^*$ and
checks all powers of the common factor $q_0$.

\subsection{The raw terminal bound}

For a retained $k$, put

\begin{equation}\label{eq:Gk}
 G_k=(S_k,m),
 \qquad S_k=\frac{z_1}{w_1}k+(\lambda-\widetilde\lambda)B.
\end{equation}

\begin{lemma}[Terminal estimate and cancellation of $q_3$]
\label{lem:terminal-cancellation}
Uniformly over every descendant in \cref{eq:Sigma8},

\begin{align}
 \Sigma_8(k)
 \ll{}&
 \frac{x^{\eps_{\rm tr}}G_k}{q_0q_3}
 \left(\frac N{\sqrt m}+\sqrt m\right)
 \left(\frac{\Delta^*}{\sqrt m}+\sqrt m\right).
 \label{eq:Sigma8-terminal}
\end{align}

Consequently

\begin{align}
 \Sigma_6(k)
 \ll{}&
 \frac{x^{3\eta}G_k}{q_0}
 \frac{\Delta_1}{\Delta^*}
 \left(\frac N{\sqrt m}+\sqrt m\right)
 \left(\frac{\Delta^*}{\sqrt m}+\sqrt m\right).
 \label{eq:Sigma6-raw}
\end{align}
\end{lemma}

\begin{proof}
Apply \cref{lem:congruence-trace} to \cref{eq:Sigma8} with
$q=q_0q_3$.  This divisor is squarefree by
\cref{lem:prime-support}.  The numerator is $w_2A_1L_1$ and
\cref{eq:L1} gives

\[
 (w_2A_1L_1,m)=(A_1A_2S_k,m)=G_k.
\]

The scale inequalities $N_2\le N$, $\Delta_2/z_1\le\Delta^*$ then give
\cref{eq:Sigma8-terminal}.

Insert this estimate into \cref{eq:Sigma7-to-8}.  Its factor $q_3$ cancels
the denominator $q_3$ exactly; this cancellation occurs before any gcd
averaging.  Then use \cref{eq:Sigma6-to-7}.  Since
$2\eta+\eps_{\rm tr}+o(1)<3\eta$ for sufficiently large $x$, the result is
\cref{eq:Sigma6-raw}.
\end{proof}

\subsection{The gcd average without
\texorpdfstring{$T\ge K$}{T >= K}}

The range \cref{eq:k-range} and the condition $w_2\mid k$ lead to
$k=w_2k_1$ and

\begin{equation}\label{eq:K1}
 |k_1|\le K_1\ll
 \frac{w_1|\Lambda|N}{x^{5\eta}w_2\Delta_1}.
\end{equation}

Because $(z_1/w_1,m)=1$,

\begin{equation}\label{eq:linear-gcd-coeff}
 \left((z_1/w_1)w_2,m\right)=(w_2,m)=t.
\end{equation}

The elementary \cref{lem:gcd-average} therefore gives

\begin{equation}\label{eq:Gk-average}
 \sum_{\substack{|k_1|\le K_1\\G_{w_2k_1}\le T}}
 G_{w_2k_1}
 \le \tau(m)(2tK_1+T).
\end{equation}

No comparison between $K_1$ and $T$ is assumed.  From
\cref{eq:Z5-ranges-c,eq:tT},

\begin{equation}\label{eq:K-over-T}
 \frac{K_1}{T}
 \ll S_x^{O(1)}q_0x^{-100\eta}
 \frac{\min\{t,H\}}{w_2}
 \le S_x^{O(1)}q_0x^{-100\eta}.
\end{equation}

Define

\begin{equation}\label{eq:Lloss}
 \mathcal L=1+K_1/T.
\end{equation}

Since $m\le x^7$, the standard divisor bound gives
$\tau(m)\ll_\eta x^\eta$.  Dividing \cref{eq:Gk-average} by a sufficiently
large constant times $x^\eta tT\mathcal L$ produces nonnegative weights
$\xi_k$ with total mass at most one.  Hence \cref{eq:Sigma6-raw} may be
written

\begin{align}
 \Sigma_6(k)\ll{}&\xi_k\mathcal L
 \frac{x^{4\eta}tT}{q_0}
 \frac{\Delta_1}{\Delta^*}
 \left(\frac N{\sqrt m}+\sqrt m\right)
 \left(\frac{\Delta^*}{\sqrt m}+\sqrt m\right).
 \label{eq:Sigma6-normalized}
\end{align}

The loss $\mathcal L$ appears once.  There is no later Cauchy--Schwarz step
that squares it.

\subsection{The genuine short scale}

The minimum in \cref{eq:Delta-star} cannot be replaced by $\Delta_1$.
Using \cref{eq:Z5-ranges-c} gives

\[
 \frac{N}{|\Lambda|x^{5\eta}}
 \gg\frac{w_1gN}{q_0x^{\delta+10\eta}H^2}.
\]

Together with \cref{eq:Delta1} and the harmless preprocessing losses this
proves, for one absolute $C_0$,

\begin{equation}\label{eq:Delta-star-lower}
 \Delta^*\gg S_x^{-C_0}
 \frac{N}{q_0^2x^{\delta+55\eta}H^2}.
\end{equation}

For the term containing both $\Delta_1$ and $\Delta^*$ we use the exact
identity

\begin{equation}\label{eq:Delta-star-branches}
 \frac1{\Delta_1\Delta^*}
 =\max\left\{
 \frac{|\Lambda|x^{5\eta}}{N\Delta_1},
 \frac1{\Delta_1^2}
 \right\}.
\end{equation}

\subsection{The exponent ledger}

Set

\begin{equation}\label{eq:Bfactor}
 \mathcal B=\min\left\{
 \frac1{q_0^2x^{\delta+10\eta}H^3},
 \frac1{H^4}
 \right\}.
\end{equation}

After substituting $T$ from \cref{eq:tT} into
\cref{eq:Sigma6-normalized}, the bound required in
\cref{eq:Sigma6-target} follows once each of

\begin{equation}\label{eq:three-terminal-terms}
 \frac{N\Delta^*}{m},\qquad N,\qquad m
\end{equation}

is at most

\begin{equation}\label{eq:terminal-RHS}
 \frac{\Delta^*}{\Delta_1}\mathcal B
 \frac{gq_0^2g_0N\Delta_1}{x^{\delta+131\eta}H^2}.
\end{equation}

Indeed, expanding the two completion factors gives

\[
 \left(\frac N{\sqrt m}+\sqrt m\right)
 \left(\frac{\Delta^*}{\sqrt m}+\sqrt m\right)
 =\frac{N\Delta^*}{m}+N+\Delta^*+m,
\]

and $\Delta^*\le N$ absorbs the third summand into the second.

After cancelling \cref{eq:terminal-RHS}, define

\begin{align}
 E_1&=\frac{x^{\delta+131\eta}}{gq_0^2g_0m}
 \max\{q_0^2x^{\delta+10\eta}H^5,H^6\},\notag\\
 E_2&=\frac{x^{\delta+131\eta}}{gq_0^2g_0\Delta^*}
 \max\{q_0^2x^{\delta+10\eta}H^5,H^6\},\notag\\
 E_3&=\frac{mx^{\delta+131\eta}}{gq_0^2g_0N\Delta^*}
 \max\{q_0^2x^{\delta+10\eta}H^5,H^6\}.
 \label{eq:E123}
\end{align}

Thus it suffices to prove $\mathcal LE_j\ll1$ for $j=1,2,3$.

From the upper bound for $\Delta_1$, the lower bound for $m$, and
$RQ^2=x^{-\eta}q_0MH$, one obtains

\begin{equation}\label{eq:m-lower-ledger}
 m\gg\frac{q_0^2x^{54\eta}MH^4}{gN}.
\end{equation}

Using $M\asymp x/N$ and
$H\ll q_0^{-1}x^{4\omega+\delta+7\eta}$, the two branches of $E_1$ are

\begin{equation}\label{eq:E1-ledger}
 E_1\ll
 x^{-1+2\gamma+4\omega+3\delta+94\eta}
 +x^{-1+2\gamma+8\omega+3\delta+91\eta}.
\end{equation}

The lower bound \cref{eq:Delta-star-lower} gives

\begin{equation}\label{eq:E2-ledger}
 E_2\ll S_x^{O(1)}\left(
 x^{28\omega+10\delta-\gamma+245\eta}
 +x^{32\omega+10\delta-\gamma+242\eta}
 \right).
\end{equation}

For $E_3$, the upper bound for $m$ first gives

\begin{equation}\label{eq:m-upper-ledger}
 m\ll\frac{q_0x^{\delta-\eta}MH^2}{g\Delta_1}.
\end{equation}

The two exact alternatives in \cref{eq:Delta-star-branches} satisfy

\begin{align}
 \frac{|\Lambda|x^{5\eta}}{N\Delta_1}
 &\ll\frac{q_0^3x^{2\delta+65\eta}H^4}{N^2},
 \label{eq:branch-Lambda}\\
 \frac1{\Delta_1^2}
 &\ll\frac{q_0^4x^{2\delta+110\eta}H^4}{N^2}.
 \label{eq:branch-Delta}
\end{align}

Substitution in $E_3$ gives four branches:

\begin{align}
 E_3\ll S_x^{O(1)}\bigl(&
 x^{1+44\omega+16\delta-4\gamma+282\eta}
 +x^{1+48\omega+16\delta-4\gamma+279\eta}\notag\\
 &+x^{1+44\omega+16\delta-4\gamma+327\eta}
 +x^{1+48\omega+16\delta-4\gamma+324\eta}
 \bigr).
 \label{eq:E3-ledger}
\end{align}

Before dropping nonpositive factors, the residual powers of $q_0$ in the
eight branches of \cref{eq:E1-ledger,eq:E2-ledger,eq:E3-ledger} are

\begin{equation}\label{eq:q0-ledger}
 -3,-6,-5,-8,-7,-10,-6,-9.
\end{equation}

The second term of $\mathcal L$ in \cref{eq:K-over-T} adds one to each
power of $q_0$ and subtracts $100\eta$ from its power of $x$.  The powers
become

\[
 -2,-5,-4,-7,-6,-9,-5,-8,
\]

and remain nonpositive.  Hence $\mathcal L$ introduces no new parameter
condition.

The $\eta$-free exponents in \cref{eq:E1-ledger,eq:E2-ledger,eq:E3-ledger}
are bounded respectively by

\begin{equation}\label{eq:three-margins}
 -1+8\omega+4\delta+2\gamma,\qquad
 32\omega+10\delta-\gamma,\qquad
 1+48\omega+16\delta-4\gamma.
\end{equation}

They are strictly negative by
\cref{eq:S52-a,eq:S52-b,eq:S52-c}.  The choice
\cref{eq:eta-choice} absorbs every displayed coefficient of $\eta$, the
fixed $S_x^{O(1)}=x^{o(1)}$ losses, and $\eps_{\rm tr}$.  Therefore
$\mathcal LE_j\ll1$ for all three $j$.

\begin{proposition}[Uniform duplicated prescribed-factor estimate]
\label{prop:duplicated-S52}
For every admissible outer localization,

\begin{equation}\label{eq:duplicated-conclusion}
 \Sigma_1^\sharp\ll g_0RQNUV^2x^{-4\eta},
\end{equation}

uniformly in all outer variables and all compatible descendant choices.
\end{proposition}

\begin{proof}
The exponent ledger proves the pointwise estimate
\cref{eq:Sigma6-target} with the normalized weights from
\cref{eq:Sigma6-normalized}.  Apply
\cref{lem:remove-ntilde,lem:diagonal-removal,lem:qvdc-supported,lem:prepared-factor}
in that order.  Uniformity of every trace and Taylor constant follows from
\cref{lem:finite-jet}.  This yields \cref{eq:duplicated-conclusion}.
\end{proof}
